\section{Results and Discussion}

Table~\ref{tab:method-comparison} summarizes the average metrics across all evaluated configurations. Method C achieves the highest precision (\(0.569\) at \(K=4\)), while Method D's GMM achieves high recall (\(0.897\) at \(K=3\), \(0.630\) at \(K=4\)). The \(K=4\) GMM configuration improves Dice (\(0.426\) to \(0.481\)) and PLL (\(-2.92\) to \(-2.77 \times 10^6\)), indicating a better probabilistic fit to the RGB distributions.

\begin{table}[t]
  \centering
  \resizebox{\linewidth}{!}{%
\begin{tabular}{@{}lcccccc@{}}
    \toprule
Method & \(K\) & Accuracy & Precision & Recall & Dice & \(\text{PLL}~(10^6)\) \\
    \midrule
Method B (LDA, supervised) & -- & -- & -- & -- & -- & -- \\
Method C (k-means, no preprocessing) & 3 & 0.927 & 0.444 & 0.925 & 0.544 & -- \\
Method C (k-means, no preprocessing) & 4 & 0.953 & 0.569 & 0.817 & 0.601 & -- \\
Method D (GMM) & 3 & 0.939 & 0.272 & 0.897 & 0.426 & -2.92 \\
Method D (GMM) & 4 & 0.951 & 0.320 & 0.630 & 0.481 & -2.77 \\
    \bottomrule
  \end{tabular}%
  }
  \caption{Cross-method comparison of segmentation metrics.}
  \label{tab:method-comparison}
\end{table}

Method B establishes a supervised baseline using LDA, exploiting strong fluorescence contrast under UV-C illumination. Method C provides a fast hard-clustering baseline, with \(K=3\) offering balanced precision and recall, while \(K=4\) achieves higher precision (\(0.569\)) but lower recall (\(0.817\)). Method D's probabilistic approach converts k-means centers into soft assignments, with PLL-based model selection identifying the best of ten EM runs. The probability maps recover gaps that hard clustering misses, achieving recall of \(0.897\) at \(K=3\) and improving Dice from \(0.426\) to \(0.481\) at \(K=4\).

Future work will expand the dataset and explore hybrid models, including supervised backbones such as U-Net or transformer encoders, while additional substrates and illumination settings will test generalization.

\section{Conclusion}

This report presents four segmentation methods for saliva detection under UV-C illumination: Method A's MATLAB thresholding, Method B's supervised LDA, Method C's k-means clustering, and Method D's likelihood-aware GMM. Method D achieves high recall while improving Dice at \(K=4\) through PLL-guided model selection, demonstrating that probabilistic modeling with multi-start EM delivers stable masks for safety-critical contamination imagery.

